% Transcription of Cayley's Collected Mathematical Papers, Volume IV
% PDF pages 451--475 (printed pages 429--453)
% Papers: 280 "On the Conics which Touch Four Given Lines"
%         281 "Note on the Wave Surface"
%         282 "On a Particular Case of Castillon's Problem"
%         283 "On a Theorem relating to Homographic Figures"
%         284 "On a New Analytical Representation of Curves in Space" (opening)

\documentclass[11pt]{article}
\usepackage[a4paper,margin=0.65in]{geometry}
\usepackage[T1]{fontenc}
\usepackage{amsmath,amssymb,amsfonts}
\usepackage{graphicx}
\usepackage{pdfpages}
\usepackage[english,shorthands=off]{babel}
\usepackage{fancyhdr}
\usepackage[bookmarks=false,colorlinks=false]{hyperref}
\usepackage[protrusion=true,expansion=false]{microtype}
\emergencystretch=3em
\tolerance=600
\providecommand{\modd}{\mathrm{mod}}
\providecommand{\Disc}{\mathrm{Disc}}
\providecommand{\canont}{}
\providecommand{\asterism}{*}
\providecommand{\Hessian}{H}
\providecommand{\tome}{}
\providecommand{\page}{p.}
\pagestyle{fancy}\fancyhf{}\fancyhead[L]{Pages 451--475}\fancyhead[R]{Cayley --- Collected Papers, Vol. IV}\renewcommand{\headrulewidth}{0.4pt}
\setlength{\headheight}{15pt}

\begin{document}

% ============================================================
% Page 429 (PDF 451) -- Paper 280 begins
% ============================================================
\begin{center}
\textsc{280]} \hfill \hfill 429
\end{center}

\begin{center}
\textbf{\large 280.}\\[1ex]
\textsc{ON THE CONICS WHICH TOUCH FOUR GIVEN LINES.}
\end{center}

\begin{center}
\textit{[From the Quarterly Journal of Pure and Applied Mathematics, vol. \textsc{iii}.\ (1860),
pp. 94--96.]}
\end{center}

\medskip

There are considerable practical difficulties in drawing a figure of the system of
conics which touch four given lines, but a notion of the figure of the system may be
obtained as follows:

Figure 20 may be taken to represent any quadrilateral whatever, having all its
sides real; and if we attend only to the unshaded spaces, it will be seen that there
are five regions which are called the inner, upper, lower, right-hand, and left-hand
regions respectively. The inner diagonals are $AC$ the vertical diagonal and $BD$ the
horizontal diagonal; the former of these traverses the inner, upper, and lower regions;
the latter the inner, right-hand, and left-hand regions; the outer diagonal is $EF$
which traverses the upper region and the right-hand and left-hand regions. The
inner or vertical and horizontal diagonals meet in the inner region, the vertical and
outer diagonals meet in the upper region, the horizontal and outer diagonals meet in

\medskip
\textit{[Figure 20: a complete quadrilateral $ABCD$ with extended sides meeting at the
diagonal points $E$, $F$; the five regions (inner, upper, lower, right-hand, left-hand)
are distinguished by shading.]}
\medskip

% ============================================================
% Page 430 (PDF 452)
% ============================================================
\begin{center}
430 \hfill ON THE CONICS WHICH TOUCH FOUR GIVEN LINES. \hfill \textsc{[280}
\end{center}

the left-hand region. No conic touching the four lines lies wholly or in part in the
shaded regions, and every conic touching the four lines lies wholly in the inner region
or wholly in the upper region, or partly in the upper and partly in the lower region,
or partly in the right-hand and partly in the left-hand region. It will be convenient
to consider, $1^{\circ}$. the conics which lie in the inner region; $2^{\circ}$. the
conics which lie in the upper and lower regions, or in the upper region only;
$3^{\circ}$. the conics which lie in the right-hand and left-hand regions.

$1^{\circ}$. The conics in the inner region are obviously ellipses; an extreme term is
the finite right line $BD$, considered as an indefinitely thin ellipse; this gradually
broadens out and there is (as a mean term) an ellipse which touches the four lines in
the points in which they are intersected two and two by the lines joining the points
$E,F$ with the point of intersection of the diagonals $AC$ and $BD$, the ellipse then
narrows in the transverse direction and at length reduces itself to the finite line
$AC$ considered as an indefinitely thin ellipse.

$2^{\circ}$. Considering first the conics which lie in the upper and lower regions,
these are of course hyperbolas, an extreme term is the infinite portions $A\infty$ and
$C\infty$ of the line $AC$, considered as an indefinitely thin hyperbola: we have then
hyperbolas such that for each of them, one branch lies in the lower region and touches
the two lines through $A$, while the other branch lies in the upper region and touches
the two lines through $C$; the points of contact of the lower branch with the lines
through $A$ gradually recede from $A$, but the point of contact on the line $AD$, which
adjoins the left-hand region, recedes with the greater rapidity, and it at last becomes
infinite while the point of contact with the line $AB$ which adjoins the right-hand
region remains finite: we have thus a hyperbola having the line $AD$ for its asymptote;
the point at infinity of this line belongs of course indifferently to the upper and
lower regions; and we may therefore consider one branch as lying in the lower region
and touching $AB$ and (at infinity) $AD$; the other branch as lying in the upper region
and touching the two lines through $C$, and besides (at infinity) the line $AD$. We
have next a series of hyperbolas such that for each of them, one branch lies in the
lower region and touches only the line $AB$, the other branch lies in the upper region
and touches the two lines through $C$ and besides the line $AD$. We arrive again at a
limiting case when the point of contact on the line $AB$ passes off to infinity, or
$AB$ becomes an asymptote; we have here in the lower region a branch touching (at
infinity) $AB$ and in the upper region a branch touching the two lines through $C$,
the line $AD$, and besides (at infinity) the line $AB$. To this succeeds a series of
hyperbolas such that for each of them, one branch lies in the lower region but does
not touch either of the lines through $A$, the other branch lies in the upper region
and touches as well the two lines through $C$ as also the two lines through $A$.
Finally the branch in the lower region passes off to infinity, or the conic becomes
a parabola lying wholly in the upper region.

We have thus arrived at the conics which lie wholly in the upper region, the
extreme term being a parabola: this passes into an ellipse touching of course the
four lines, and gradually reducing itself to the finite line $EF$ considered as an
indefinitely thin ellipse.

% ============================================================
% Page 431 (PDF 453)
% ============================================================
\begin{center}
\textsc{280]} \hfill ON THE CONICS WHICH TOUCH FOUR GIVEN LINES. \hfill 431
\end{center}

$3^{\circ}$. We come now to the conics which lie in the right-hand and left-hand
regions; an extreme term is the conic composed of the infinite portions $E\infty$ and
$F\infty$ of the line $EF$, considered as an indefinitely thin hyperbola; we have then
a series of hyperbolas such that for each of them the branch in the right-hand region
touches the two lines through $E$, and the branch in the left-hand region the two
lines through $F$; we then arrive at a limiting case where the branch in the
right-hand region touches the line $EC$ the upper boundary of the right-hand region
at $\infty$, or where $EC$ is an asymptote; we may say that the branch in the
right-hand region touches the line $EB$ and (at infinity) the line $EC$, while the
branch in the left-hand region touches not only the two lines through $F$ but also
(at infinity) the line $EC$; we have next a series of hyperbolas such that for each
of them the branch in the right-hand region touches only the line $EB$ while the
branch in the left-hand region touches as well the two lines through $F$ as the line
$EC$; we have then again a limiting case, viz.\ the branch in the right-hand region
touches the line $BC$ at infinity, or $BC$ is an asymptote; we may here say that the
branch in the right-hand region touches the line $BE$ and (at infinity) the line $BC$,
while the branch in the left-hand region touches the lines through $D$ and (at
infinity) the line $BC$; and to this succeeds a series of hyperbolas such that for
each of them the branch in the right-hand region touches the two lines through $B$
while the branch in the left-hand region touches the two lines through $D$; the
hyperbola finally becomes the infinite portions $B\infty$ and $D\infty$ of the line
$BD$, considered as an indefinitely thin hyperbola.

It is to be noticed that the entire system commences with the finite line $BD$
considered as an indefinitely thin ellipse and passes on continuously to the infinite
portions $B\infty$ and $D\infty$ of the same line, considered as an indefinitely thin
hyperbola: it is therefore a cyclical series, and any term whatever might have been
taken for the commencement of the series.

\bigskip
% ============================================================
% Page 432 (PDF 454) -- Paper 281 begins
% ============================================================
\begin{center}
432 \hfill \hfill \textsc{[281}
\end{center}

\begin{center}
\textbf{\large 281.}\\[1ex]
\textsc{NOTE ON THE WAVE SURFACE.}
\end{center}

\begin{center}
\textit{[From the Quarterly Journal of Pure and Applied Mathematics, vol. \textsc{iii}.\ (1860),
pp. 142--144.]}
\end{center}

\medskip

In the paper ``On the Wave Surface'' in the last Number of the \textit{Journal},
[277], I stated that it was shown in the second of Zech's Memoirs that the curves of
contact of the developables $F$ and $G$ with the wave surface were the curves of
curvature of the wave surface. I had not examined the demonstration of this theorem,
and I had overlooked the author's note ``Die Kr\"ummungslinien der Wellenfl\"ache
zweiaxiger Krystalle, \&c.,'' \textit{Crelle}, t.\ \textsc{liv}.\ p.\ 94 (Feb.\ 1858),
where he points out that in a phrase which he quotes, he had assumed without
demonstration a theorem which was in fact erroneous, and that he retracted all that
he had said in No.\ 11 of his Memoir (the portion which contains the theorem as to
the curves of curvature of the wave surface). M.\ Bertrand in a note in the
\textit{Comptes Rendus}, t.\ \textsc{xlvii}.\ pp.\ 817--819 (Nov.\ 1858), after
referring to my paper, remarks that the theorem as to the curves of curvature of the
wave surface appeared to him so remarkable that he hastened to investigate a proof
of it, but that he very soon discovered that the theorem was unfortunately erroneous;
and he proceeds to show why the theorem cannot be true. M.\ Bertrand's demonstration
is as follows:

\textsc{Theorem} I. If from a point $O$, we let fall perpendiculars on the tangent
planes of a surface, the locus of their feet is a new surface. Let $P$ be a point of
this surface corresponding to the point $M$ of the first surface, the normal at $P$
passes through the middle point of $OM$.

\textsc{Theorem} II. If the curve of curvature of a surface is such that the tangent
planes at the several points of the curve are equidistant from a point $O$, the curve
of curvature is situate on a sphere having the point $O$ for its centre.

For suppose that we let fall from the point $O$ perpendiculars on the tangent
planes to the surface at the several points of the curve of curvature in question. The

% ============================================================
% Page 433 (PDF 455)
% ============================================================
\begin{center}
\textsc{281]} \hfill NOTE ON THE WAVE SURFACE. \hfill 433
\end{center}

locus of the feet will be a spherical curve, a normal at any point $P$ of the curve
will it is clear pass through the point $O$; besides, in virtue of the first theorem,
another normal will pass through the middle point of the radius $OM$ drawn to the
corresponding point $M$ of the surface; the tangent to the curve which is the locus of
the points $P$ is therefore perpendicular to the plane $MOP$, and consequently to the
line $MP$. But when a developable surface is circumscribed about a sphere, the
perpendiculars let fall from the centre of the sphere on the tangent planes of the
developable surface have their feet on the generating lines; and consequently the
curve which is the locus of the points $P$ is situate on the developable surface
(viz., the developable surface enveloped by the tangent planes at the several points
of the curve of curvature of the given surface) and cuts the generating lines at
right angles. But the curve, the locus of the point $M$, being by hypothesis a curve
of curvature of the given surface, also cuts at right angles the generating lines of
the developable surface, the two curves are therefore equidistant curves on the
developable surface, viz.\ $MP$ is constant, and since by hypothesis $OP$ is constant,
$OM$ is also constant, which proves the second theorem.

This being premised, it is to be recollected that the wave surface may be generated
in two different ways; $1^{\circ}$. it is the locus of the extremities of the central
perpendiculars to the diametral sections of an ellipsoid $E$, equal to the axes of
these sections; $2^{\circ}$. it is the envelope of planes parallel to the diametral
sections of a second ellipsoid $E'$ at distances inversely proportional to the axes
of the sections. To obtain all the tangent planes of the wave surface situate at a
distance $h$ from the centre, it is necessary to find in the ellipsoid the diametral
sections which have an axis equal to $\dfrac{1}{h}$; for this, we may cut the ellipsoid
by a concentric sphere of the radius $\dfrac{1}{h}$, and draw tangent planes to the
cone having for its vertex the centre of the ellipsoid and passing through the curve
of intersection with the sphere. The tangent planes of the wave surface respectively
parallel to the tangent planes of the cone will be at the distance $h$ from the
centre, and, if they touched the wave surface along a curve of curvature, it would
follow from the foregoing theorem II.\ that their points of contact would be all at
the same distance from the centre; but the distance of the centre of the wave surface
from the point of contact of a tangent plane is inversely proportional to the
perpendicular let fall from the centre of the ellipsoid $E'$ on the tangent plane at
the corresponding point of this ellipsoid: it follows that the curve of intersection
of an ellipsoid with a concentric sphere is not such that the tangent planes of the
ellipsoid at the several points of the curve are at the same distance from the centre,
and consequently the tangent planes which were under consideration do not determine
by their points of contact a curve of curvature of the wave surface.

I remark in relation to M.\ Bertrand's theorem I.\ that it is immediately connected
with a well-known theorem which occurs in optics---viz., if rays proceeding from a
point are reflected at a surface, and if from the radiant point to the surface and
thence along the reflected ray, we measure off a constant distance, the surface which
is the locus of the point so obtained (the secondary focal surface of Dandelin and
Quetelet)

\begin{center}
C.\ IV. \hfill 55
\end{center}

% ============================================================
% Page 434 (PDF 456)
% ============================================================
\begin{center}
434 \hfill NOTE ON THE WAVE SURFACE. \hfill \textsc{[281}
\end{center}

is an orthogonal trajectory of the reflected rays. In fact, if $O$ be the radiant
point, and $OM'$ the incident ray, and if from $M'$ we measure off on the reflected
ray a distance $=-OM'$, that is, on the reflected ray produced backwards, a distance
$M'P=OM'$, then the whole distance from $O$ is $OM'-OM'=0$, and the surface which is
the locus of the point $P$ is consequently an orthogonal trajectory to the reflected
rays. But the point $P$ may, it is clear, be constructed as follows: viz., on the
tangent plane at $M'$ let fall the perpendicular $OP'$, and produce it to a point $P$
such that $OP'=P'P$. And if we produce $OM'$ to $M$ so that $OM'=M'M$, then it is
clear that the locus of $M$ is a surface similar and similarly situated with the
original surface, but of double the magnitude, and that $OP$ is the perpendicular
from $O$ upon the tangent plane at $M$ of the last-mentioned surface. And, by what
precedes, the line $PM'$ from $P$ to the middle point of $OM$ is a normal of the
surface which is the locus of $P$: the corresponding theorem \textit{in plano} is in
fact actually given by Dandelin.

\medskip
\hfill 31\textit{st Dec.}, 1858.

\bigskip
% ============================================================
% Page 435 (PDF 457) -- Paper 282 begins
% ============================================================
\begin{center}
\textsc{282]} \hfill \hfill 435
\end{center}

\begin{center}
\textbf{\large 282.}\\[1ex]
\textsc{ON A PARTICULAR CASE OF CASTILLON'S PROBLEM.}
\end{center}

\begin{center}
\textit{[From the Quarterly Journal of Pure and Applied Mathematics, vol. \textsc{iii}.\ (1860),
pp. 157--164.]}
\end{center}

\medskip

The problem referred to as Castillon's problem, being itself a particular case of
the problem of the in-and-circumscribed triangle, is as follows: viz.\ ``In a circle
to inscribe a triangle the sides of which pass through three given points,'' and the
reciprocal problem of course presents itself ``about a circle to circumscribe a
triangle, the angles of which lie in given lines.'' If in Castillon's problem the
three given points are the angles of a triangle circumscribed about the circle, or
if in the reciprocal problem the three given lines are the sides of a triangle
circumscribed about the circle, we have a circle and an inscribed and circumscribed
triangle, such that the sides of the inscribed triangle pass through the angles of
the circumscribed triangle, and the problem arises ``given one of the triangles to
determine the other triangle.'' The problem, so far as I am aware, was first proposed
by Clausen, who has given, \textit{Crelle}, t.\ \textsc{iv}.\ (1829), p.\ 391, a very
elegant solution, which I propose to reproduce here.

Let the angle of a point be defined as the inclination to a fixed radius, of the
line from the centre through the given point.

Let $\alpha,\beta,\gamma$ be the angles of the points of contact of the sides of the
circumscribed triangle.

$\alpha',\beta',\gamma'$, the angles of the angular points of the circumscribed triangle.

$a,b,c$, the angles of the angular points of the inscribed triangle.

$a',b',c'$, the angles of the intersections of the perpendiculars from the centre
on the sides of the inscribed triangle.

Therefore
\[ 2a'=b+c,\qquad 2\alpha'=\beta+\gamma, \]
\[ 2b'=c+a,\qquad 2\beta'=\gamma+\alpha, \]
\[ 2c'=a+b,\qquad 2\gamma'=\alpha+\beta. \]

\begin{center}55--2\end{center}

% ============================================================
% Page 436 (PDF 458)
% ============================================================
\begin{center}
436 \hfill ON A PARTICULAR CASE OF CASTILLON'S PROBLEM. \hfill \textsc{[282}
\end{center}

Hence observing that the distance of one of the angles of the circumscribed triangle
is $\sec\tfrac{1}{2}(\beta-\gamma)$, and that the projection of this line perpendicular
to the corresponding side of the inscribed triangle is equal to
$\sec\tfrac{1}{2}(\beta-\gamma)\cos(\alpha'-a')$, which is equal to the projection
of the radius perpendicular to the same side, or to $\cos\tfrac{1}{2}(b-c)$, we have
\[ \cos(\alpha'-a')=\cos\tfrac{1}{2}(\beta-\gamma)\cos\tfrac{1}{2}(b-c), \]
\[ \cos(\beta'-b')=\cos\tfrac{1}{2}(\gamma-\alpha)\cos\tfrac{1}{2}(c-a), \]
\[ \cos(\gamma'-c')=\cos\tfrac{1}{2}(\alpha-\beta)\cos\tfrac{1}{2}(a-b). \]

Write
\[ b-c=4f,\qquad b+c-2\alpha=4x, \]
\[ c-a=4g,\qquad c+a-2\beta=4y, \]
\[ a-b=4h,\qquad a+b-2\gamma=4z; \]
therefore
\[ \cos(y+z+g-h)=\cos 2f\cos(f+y-z), \]
\[ \cos(z+x+h-f)=\cos 2g\cos(g+z-x), \]
\[ \cos(x+y+f-g)=\cos 2h\cos(h+x-y), \]
or, since $f+g+h=0$,
\[ \cos\{(y-h)+(z+g)\}=\cos 2f\cos\{(y-h)-(z+g)\}, \]
\[ \cos\{(z-f)+(x+h)\}=\cos 2g\cos\{(z-f)-(x+h)\}, \]
\[ \cos\{(x-g)+(y+f)\}=\cos 2h\cos\{(x-g)-(y+f)\}, \]
or
\[ \tan^{2}f=\tan(y-h)\tan(z+g), \]
\[ \tan^{2}g=\tan(z-f)\tan(x+h), \]
\[ \tan^{2}h=\tan(x-g)\tan(y+f). \]
Put
\[ y-h=\zeta,\quad z+g=\zeta_{1}, \]
\[ z-f=\eta,\quad x+h=\eta_{1}, \]
\[ x-g=\xi,\quad y+f=\xi_{1}, \]
\[ \tan f=l,\quad \tan\xi=x,\quad \tan\xi_{1}=x_{1}, \]
\[ \tan g=m,\quad \tan\eta=y,\quad \tan\eta_{1}=y_{1}, \]
\[ \tan h=n,\quad \tan\zeta=z,\quad \tan\zeta_{1}=z_{1}. \]
We have then
\[ \tan^{2}f=\tan\zeta\tan\zeta_{1}, \]
\[ \tan^{2}g=\tan\eta\tan\eta_{1}, \]
\[ \tan^{2}h=\tan\xi\tan\xi_{1}, \]
or, what is the same thing,
\[ l^{2}=xx_{1},\qquad m^{2}=yy_{1},\qquad n^{2}=zz_{1}. \]

% ============================================================
% Page 437 (PDF 459)
% ============================================================
\begin{center}
\textsc{282]} \hfill ON A PARTICULAR CASE OF CASTILLON'S PROBLEM. \hfill 437
\end{center}

But to obtain equations involving only one of the sets $(x,y,z),(x_{1},y_{1},z_{1})$,
it is proper to write
\[ \tan^{2}f=\tan(\zeta+g)\tan\zeta=\tan\zeta_{1}\tan(\zeta_{1}-h), \]
\[ \tan^{2}g=\tan(\eta+h)\tan\eta=\tan\eta_{1}\tan(\eta_{1}-f), \]
\[ \tan^{2}h=\tan(\xi+f)\tan\xi=\tan\xi_{1}\tan(\xi_{1}-g); \]
taking the first set of equations, we have
\[ l^{2}=\dfrac{x(z+m)}{1-mz}\cdot\dfrac{z(1-nx)}{x+n}, \]
therefore
\[ z=\dfrac{l^{2}-mx}{x+l^{2}n}=\dfrac{m^{2}(1-nx)}{x+n}; \]
therefore
\[ y+l\,:\,1-ly=nl^{2}+ln+(1-m^{2}n)x\,:\,n-lm^{2}+(1+lm^{2}n)x; \]
therefore
\[ m^{2}\{n-lm^{2}+(1+lm^{2}n)x\}(l^{2}n+x)=(l^{2}-mx)\{m^{2}+ln+(1-m^{2}n)x\}; \]
\[ (l^{4}m^{2}+l^{2}n-l^{2}m^{4}n^{2}+l^{2}m^{4}n^{2}) \]
\[ +x(-n^{2}+lm^{2}n^{2}-l^{2}mn^{2}-l^{2}m^{4}-m^{4}-lmn+l^{2}-l^{2}m^{2}n) \]
\[ +x^{2}(-lm+m^{4}n-n^{3}-lm^{2}n^{2})=0. \]
Now $l+m+n=lmn$, and by means of this relation we find
\[ \tfrac{1}{l^{2}}\operatorname{Coef.}x^{0}=l\{lm^{2}+l^{2}n-n^{2}(l+m+n)+m(l+m+n)^{2}\}, \]
\[ \operatorname{Coef.}x=-n^{2}+mn(l+m+n)-ln(l+m+n)-(l+m+n)^{2}-m^{2}\!-lmn+l^{2}-l^{2}m^{2}n(l+m+n), \]
\[ \tfrac{1}{m^{2}}\operatorname{Coef.}x^{2}=\{-l^{2}m+m^{2}(l+m+n)-ln^{2}-n(l+m+n)^{2}\}; \]
or, reducing,
\[ \tfrac{1}{l^{2}}\operatorname{coef.}\,x^{0}=(m^{3}+2m^{2}n-n^{3})+l(3m^{2}+2mn-n^{2})+l^{2}(m+n) \]
\[ =(m+n)\{(m^{2}+mn-n^{2})+l(3m-n)+l^{2}\} \]
\[ =(m+n)\{(l+m)^{2}+(l+m)(l+n)-(l+n)^{2}\}, \]
\[ -\operatorname{coef.}\,x=m^{3}+m^{2}n+mn^{2}+n^{3}+2l(m^{2}+2mn+n^{2})+2l^{2}(m+n) \]
\[ =(m+n)\{m^{2}+n^{2}+2l(m+n)+2l^{2}\} \]
\[ =(m+n)\{(l+m)^{2}+(l+n)^{2}\}, \]
\[ \tfrac{1}{m^{2}}\operatorname{coef.}\,x^{2}=(m^{3}-2m^{2}n-n^{3})+l(n^{2}-2mn-3n^{2})-l^{2}(m+n) \]
\[ =(m+n)\{(m^{2}-mn-n^{2})+l(m-3n)-l^{2}\} \]
\[ =(m+n)\{(l+m)^{2}-(l+m)(l+n)-(l+n)^{2}\}; \]

% ============================================================
% Page 438 (PDF 460)
% ============================================================
\begin{center}
438 \hfill ON A PARTICULAR CASE OF CASTILLON'S PROBLEM. \hfill \textsc{[282}
\end{center}

and the equation becomes
\[ l^{2}\{(l+m)^{2}+(l+m)(l+n)-(l+n)^{2}\} \]
\[ -2lx\{(l+m)^{2}+(l+n)^{2}\} \]
\[ +x^{2}\{(l+m)^{2}-(l+m)(l+n)-(l+n)^{2}\}=0, \]
which may be written
\[ \resizebox{\textwidth}{!}{$\displaystyle
[l\{(l+m)^{2}-(l+m)(l+n)-(l+n)^{2}\}-x\{(l+m)^{2}+(l+n)^{2}\}]^{2}
=5(l+m)^{2}(l+n)^{2},
$} \]
and we have therefore
\[ \dfrac{x}{l}=\dfrac{(l+m)^{2}+\sqrt{(5)}(l+m)(l+n)+(l+n)^{2}}{(l+m)^{2}-(l+m)(l+n)-(l+n)^{2}}, \]
or, reducing and observing also that $l^{2}=xx_{1}$, we have
\[ \dfrac{x}{l}=\dfrac{x_{1}}{l}=\dfrac{l+m+\tfrac{1}{2}\{1+\sqrt{(5)}\}(l+n)}{(l+m)-\tfrac{1}{2}\{1+\sqrt{(5)}\}(l+n)}. \]
The values of $(l,m,n)$ are
\[ l=\tan\tfrac{1}{2}(b-c),\quad m=\tan\tfrac{1}{2}(c-a),\quad n=\tan\tfrac{1}{2}(a-b), \]
and those of $(x,y,z),(x_{1},y_{1},z_{1})$ are
\[ x=\tan\tfrac{1}{2}(b+c-2\gamma),\quad y=\tan\tfrac{1}{2}(c+a-2\alpha),\quad z=\tan\tfrac{1}{2}(a+b-2\beta), \]
\[ x_{1}=\tan\tfrac{1}{2}(b+c-2\beta),\quad y_{1}=\tan\tfrac{1}{2}(c+a-2\beta),\quad z_{1}=\tan\tfrac{1}{2}(a+b-2\alpha), \]
and the foregoing result shows that the values of
\[ \dfrac{x}{l}\ \text{or}\ \dfrac{x_{1}}{l};\quad \dfrac{y}{m}\ \text{or}\ \dfrac{y_{1}}{m};\quad \dfrac{z}{n}\ \text{or}\ \dfrac{z_{1}}{n} \]
are
\[ \dfrac{\sin\tfrac{1}{2}(a-b)+\tfrac{1}{2}\{1+\sqrt{(5)}\}\sin\tfrac{1}{2}(c-a)}{\sin\tfrac{1}{2}(a-b)-\tfrac{1}{2}\{1+\sqrt{(5)}\}\sin\tfrac{1}{2}(c-a)}, \]
\[ \dfrac{\sin\tfrac{1}{2}(b-c)+\tfrac{1}{2}\{1+\sqrt{(5)}\}\sin\tfrac{1}{2}(a-b)}{\sin\tfrac{1}{2}(b-c)-\tfrac{1}{2}\{1+\sqrt{(5)}\}\sin\tfrac{1}{2}(a-b)}, \]
\[ \dfrac{\sin\tfrac{1}{2}(c-a)+\tfrac{1}{2}\{1+\sqrt{(5)}\}\sin\tfrac{1}{2}(b-c)}{\sin\tfrac{1}{2}(c-a)-\tfrac{1}{2}\{1+\sqrt{(5)}\}\sin\tfrac{1}{2}(b-c)}. \]
Write for a moment $\dfrac{x}{l}=\dfrac{x_{1}}{l}=k$; therefore
\[ \dfrac{k-1}{1+k}\dfrac{1}{xx_{1}}\cdot\dfrac{l(k-\tfrac{1}{k})}{1+k}=-l\dfrac{P-Q}{P+Q}\ \text{if}\ k=\dfrac{P}{Q}; \]
and taking $P,Q$ for the numerator and the denominator respectively of the foregoing
fractional expression for $k$, we find

% ============================================================
% Page 439 (PDF 461)
% ============================================================
\begin{center}
\textsc{282]} \hfill ON A PARTICULAR CASE OF CASTILLON'S PROBLEM. \hfill 439
\end{center}

also
\[ P^{2}-Q^{2}=2\{1+\sqrt{(5)}\}\sin\tfrac{1}{2}(a-b)\sin\tfrac{1}{2}(c-a), \]
\[ 2PQ=2[\sin^{2}\tfrac{1}{2}(a-b)-\tfrac{1}{4}\{1+\sqrt{(5)}\}^{2}\sin^{2}\tfrac{1}{2}(c-a)] \]
\[ =-\{1+\sqrt{(5)}\}[\tfrac{1}{2}\{1-\sqrt{(5)}\}\sin^{2}\tfrac{1}{2}(a-b)+\tfrac{1}{2}\{1+\sqrt{(5)}\}\sin^{2}\tfrac{1}{2}(c-a)]; \]
also
\[ \dfrac{z-x}{1+xx_{1}}=\tan\tfrac{1}{2}(\beta-\gamma), \]
we have therefore
\[ \tan\tfrac{1}{2}(\beta-\gamma)=\dfrac{-2\sin\tfrac{1}{2}(b-c)\sin\tfrac{1}{2}(c-a)\sin\tfrac{1}{2}(a-b)}{\tfrac{1}{2}\{1-\sqrt{(5)}\}\sin^{2}\tfrac{1}{2}(a-b)+\tfrac{1}{2}\{1+\sqrt{(5)}\}\sin^{2}\tfrac{1}{2}(c-a)}, \]
\[ \tan\tfrac{1}{2}(\gamma-\alpha)=\dfrac{-2\sin\tfrac{1}{2}(b-c)\sin\tfrac{1}{2}(c-a)\sin\tfrac{1}{2}(a-b)}{\tfrac{1}{2}\{1-\sqrt{(5)}\}\sin^{2}\tfrac{1}{2}(b-c)+\tfrac{1}{2}\{1+\sqrt{(5)}\}\sin^{2}\tfrac{1}{2}(a-b)}, \]
\[ \tan\tfrac{1}{2}(\alpha-\beta)=\dfrac{-2\sin\tfrac{1}{2}(b-c)\sin\tfrac{1}{2}(c-a)\sin\tfrac{1}{2}(a-b)}{\tfrac{1}{2}\{1-\sqrt{(5)}\}\sin^{2}\tfrac{1}{2}(c-a)+\tfrac{1}{2}\{1+\sqrt{(5)}\}\sin^{2}\tfrac{1}{2}(b-c)}, \]
equations which determine the circumscribed triangle when the inscribed triangle is
given.

The more general problem, for a conic and an inscribed and circumscribed triangle
such that the sides of the inscribed triangle pass through the angles of the
circumscribed triangle, ``given one of the triangles to determine the other'' is
solved by M\"obius, \textit{Crelle}, t.\ \textsc{v}.\ (1830), p.\ 103, by means of his
Barycentric Calculus, which is in fact the method of trilinear coordinates. The
solution is in effect as follows:

Let $\xi=0$, $\eta=0$, $\zeta=0$ be the equations of the sides of the inscribed
triangle, $\xi,\eta,\zeta$ may be considered as containing each of them an implicit
constant factor, and the equation of the conic may be taken to be
\[ \eta\zeta+\zeta\xi+\xi\eta=0; \]
moreover, if $x=0$, $y=0$, $z=0$ be the equations of the sides of the circumscribed
triangle, then considering $x,y,z$ as also containing each of them an implicit
constant factor, the equation of the conic may be taken to be
\[ x^{2}+y^{2}+z^{2}-2yz-2zx-2xy=0. \]
Suppose now the sides of the inscribed triangle pass through the angles of the
circumscribed triangle, we have equations such as
\[ \xi=b'y+cz,\quad \eta=c'z+ax,\quad \zeta=a'x+by; \]
substituting these values we must have identically
\[ \eta\zeta+\zeta\xi+\xi\eta+m(x^{2}+y^{2}+z^{2}-2yz-2zx-2xy)=0, \]
\[ aa'+m=0,\quad bc+bc'+b'c=2m, \]
\[ bb'+m=0,\quad ca+ca'+c'a=2m, \]
\[ cc'+m=0,\quad ab+ab'+a'b=2m, \]

% ============================================================
% Page 440 (PDF 462)
% ============================================================
\begin{center}
440 \hfill ON A PARTICULAR CASE OF CASTILLON'S PROBLEM. \hfill \textsc{[282}
\end{center}

and, substituting for $a',b',c'$ their values,
\[ (m-bc)^{2}=mb^{2},\quad (m-ca)^{2}=mc^{2},\quad (m-ab)^{2}=ma^{2}, \]
or, if $m=n^{2}$,
\[ n^{2}-bc=nb,\quad n^{2}-ca=nc,\quad n^{2}-ab=na, \]
[the signs on the second side must be all $+$ or all $-$ and the $-$ sign may be
omitted without loss of generality].

Hence
\[ a=b=c=-\tfrac{1}{2}\{1+\sqrt{(5)}\}n,\quad \sqrt{(5)}\ \text{being written for}\ +\sqrt{(5)}. \]
\[ a'=b'=c'=-\tfrac{1}{2}\{1-\sqrt{(5)}\}n=-va, \]
if for shortness
\[ v=\dfrac{1-\sqrt{(5)}}{1+\sqrt{(5)}},\ \text{i.e.}\ v=\dfrac{\sqrt{(5)}-1}{\sqrt{(5)}+1},\ \text{or}\ v^{2}-3v+1=0, \]
and $v$ having this value, the equations give
\[ \xi=vy-z,\quad \eta=vz-x,\quad \zeta=vx-y, \]
whence also
\[ 4(2v-1)x=(3v-1)\xi+v\eta+\zeta, \]
\[ 4(2v-1)y=\xi+(3v-1)\eta+v\zeta, \]
\[ 4(2v-1)z=v\xi+\eta+(3v-1)\zeta. \]
Each side of the circumscribed triangle has on it four points, viz.\ two angles of the
circumscribed triangle, a point of contact with the conic, and an intersection with
the corresponding side of the inscribed triangle. Thus for the side $x=0$, the four
points are given as the intersections of $x=0$, with
\[ y=0,\quad z=0,\quad y-z=0,\quad vy+z=0, \]
and the anharmonic ratio of the four points is therefore a given quantity.\textsuperscript{*}

Again, each side of the inscribed triangle has on it four points, viz.\ two angles
of the inscribed triangle, a point of intersection with the tangent at the opposite
angle of the inscribed conic, and a point of intersection with the corresponding side
of the circumscribed triangle.

Thus for the side $\xi=0$, the four points are given as the points of intersection
of $\xi=0$ with the lines
\[ \eta=0,\quad \zeta=0,\quad \eta+\zeta=0,\quad \eta+(3v-1)\zeta=0, \]
and the anharmonic ratio of the four points is therefore a given quantity.\textsuperscript{*}

If we draw tangents at the angles of the inscribed triangle, we have a new triangle,
the sides of which are $\eta+\zeta=0$, $\zeta+\xi=0$, $\xi+\eta=0$, and joining the
angles of this triangle with the points of contact of the opposite sides (i.e.\ the
angles of the

% ============================================================
% Page 441 (PDF 463)
% ============================================================
\begin{center}
\textsc{282]} \hfill ON A PARTICULAR CASE OF CASTILLON'S PROBLEM. \hfill 441
\end{center}

inscribed triangle), we have three lines $\eta-\zeta=0$, $\zeta-\xi=0$, $\xi-\eta=0$
meeting in a point $\xi=\eta=\zeta$, which is obviously the same as the point
$x=y=z$, which is the point of intersection of the lines joining the angles of the
circumscribed triangle with the points of contact of the opposite sides.\textsuperscript{*}

The coordinates of the points of contact of the sides of the circumscribed triangle
are given by $(x=0,\ y-z=0)$, $(y=0,\ z-x=0)$, $(z=0,\ x-y=0)$, these points form
therefore an inscribed triangle the sides of which are
\[ y+z-x=0,\quad z+x-y=0,\quad x+y-z=0. \]

Again, the tangents at the angles of the inscribed triangle form a circumscribed
triangle the sides of which are
\[ \eta+\zeta=0,\quad \zeta+\xi=0,\quad \xi+\eta=0; \]
therefore
\[ (\xi+\eta)-v(\zeta+\eta)=(1-v)\eta-v\zeta+\xi=2vx-(1-v+v^{2})y-(1-v+v^{2})z=-2v(y+z-x), \]
and we thus have
\[ \xi+\eta-v(\zeta+\eta)=-2v(y+z-x), \]
\[ \eta+\zeta-v(\eta+\xi)=-2v(z+x-y), \]
\[ \zeta+\xi-v(\xi+\eta)=-2v(x+y-z), \]
equations which show that the sides of the second inscribed triangle pass through the
angles of the second circumscribed triangle, and that the two systems are consequently
reciprocal.\textsuperscript{*}

The four theorems marked \textsuperscript{*} are all of them contained in the paper
by M\"obius.

\medskip
\begin{center}C.\ IV.\hfill 56\end{center}

% ============================================================
% Page 442 (PDF 464) -- Paper 283 begins
% ============================================================
\begin{center}
442 \hfill \hfill \textsc{[283}
\end{center}

\begin{center}
\textbf{\large 283.}\\[1ex]
\textsc{ON A THEOREM RELATING TO HOMOGRAPHIC FIGURES.}
\end{center}

\begin{center}
\textit{[From the Quarterly Journal of Pure and Applied Mathematics, vol. \textsc{iii}.\ (1860),
pp. 177--180.]}
\end{center}

\medskip

The following theorem is not new, but I do not remember where to find it:

Given any two homographic figures; there exists in the first figure a point $S$,
which is the focus of an (ellipse or hyperbola, say an) ellipse $\Sigma$; and in the
second figure a point $S'$, which is the focus of a (hyperbola or ellipse, say a)
hyperbola $\Sigma'$, such that the points $S,S'$ and the conics $\Sigma,\Sigma'$
correspond to each other, and the conics $\Sigma,\Sigma'$ are so related that the
foci and vertices of the one may be superimposed upon the vertices and foci of the
other. Moreover the perpendiculars from $S,S'$ upon corresponding tangents of
$\Sigma,\Sigma'$ will be equal.

Write
\[ x\,:\,y\,:\,1=\alpha\xi+\beta\eta+\gamma\,:\,\alpha'\xi+\beta'\eta+\gamma'\,:\,\alpha''\xi+\beta''\eta+\gamma'', \]
then $l,m,\lambda,\mu$ may be so determined that $(x-l)+i(y-m)$ shall vanish with
$(\xi-\lambda)+i(\eta-\mu)$, for if we write
\[ J=\alpha-l\alpha'',\quad J'=\alpha'-m\alpha'', \]
\[ K=\beta-l\beta'',\quad K'=\beta'-m\beta'', \]
\[ L=\gamma-l\gamma'',\quad L'=\gamma'-m\gamma'', \]
we have
\[ (x-l)+i(y-m)=\dfrac{(J+iJ')\xi+(K+iK')\eta+(L+iL')}{\alpha''\xi+\beta''\eta+\gamma''}, \]
and therefore
\[ K+iK'=i(J+iJ'), \]
\[ L+iL'=(J+iJ')(\lambda+i\mu), \]

% ============================================================
% Page 443 (PDF 465)
% ============================================================
\begin{center}
\textsc{283]} \hfill ON A THEOREM RELATING TO HOMOGRAPHIC FIGURES. \hfill 443
\end{center}

that is
\[ K=-J',\quad K'=J, \]
or
\[ l\beta''+m\alpha''=\alpha'+\beta, \]
\[ -l\alpha''+m\beta''=-\alpha+\beta', \]
whence also
\[ l(\alpha''^{2}+\beta''^{2})=\alpha'\beta''-\alpha''\beta'+\alpha\alpha''+\beta\beta'', \]
\[ m(\alpha''^{2}+\beta''^{2})=\alpha''\beta-\alpha\beta''+\alpha\alpha'+\beta'\beta'', \]
which determine $l,m$; and then $\lambda,\mu$ are given by
\[ \lambda+i\mu=\dfrac{L+iL'}{J+iJ'}, \]
but more simply as follows, viz.\ remarking that if $c=\beta'\gamma''-\beta''\gamma'$,
\&c., so that
\[ \xi\,:\,\eta\,:\,1=ax+a'y+a''\,:\,bx+b'y+b''\,:\,cx+c'y+c'', \]
then we have
\[ \lambda(c^{2}+c'^{2})=bc'-b'c+ac+a'c', \]
\[ \mu(c^{2}+c'^{2})=ca'-c'a+bc+b'c'; \]
the equations $x=l$, $y=m$ give the point $S$ and $\xi=\lambda$, $\eta=\mu$ the
point $S'$. The values of $J,J'$ are
\[ J=\dfrac{-c\alpha''-c'\beta''}{\alpha''^{2}+\beta''^{2}},\quad J'=\dfrac{-c'\alpha''+c\beta''}{\alpha''^{2}+\beta''^{2}}, \]
and therefore
\[ J^{2}+J'^{2}=\dfrac{c^{2}+c'^{2}}{\alpha''^{2}+\beta''^{2}}=\dfrac{k^{2}}{K^{2}}; \]
if
\[ k=\sqrt{(c^{2}+c'^{2})},\qquad K=\sqrt{(\alpha''^{2}+\beta''^{2})}; \]
we have therefore
\[ k=\sqrt{(c^{2}+c'^{2})},\quad K=\sqrt{(\alpha''^{2}+\beta''^{2})}. \]

Consider now the expression
\[ (x-l)\cos\vartheta+(y-m)\sin\vartheta-\varpi, \]
which made equal to 0 would be the equation of a line the perpendicular distance
of which from $S$ is $\varpi$, and which distance is inclined at an angle $\vartheta$
to the axis of $x$; then
\begin{align*}
(x-l)\cos\vartheta+(y-m)\sin\vartheta
 &= \tfrac{1}{2}[(\cos\vartheta-i\sin\vartheta)\{(x-l)+i(y-m)\}\\
 &\qquad +(\cos\vartheta+i\sin\vartheta)\{(x-l)-i(y-m)\}]\\
 &= \tfrac{1}{2(\alpha''\xi+\beta''\eta+\gamma'')}\bigl[(\cos\vartheta-i\sin\vartheta)\{(\xi-\lambda)+i(\eta-\mu)\}(J+iJ')\\
 &\qquad +(\cos\vartheta+i\sin\vartheta)\{(\xi-\lambda)-i(\eta-\mu)\}(J-iJ')\bigr]
\end{align*}

\begin{center}56--2\end{center}

% ============================================================
% Page 444 (PDF 466)
% ============================================================
\begin{center}
444 \hfill ON A THEOREM RELATING TO HOMOGRAPHIC FIGURES. \hfill \textsc{[283}
\end{center}

\[ \resizebox{\textwidth}{!}{$\displaystyle
=\dfrac{k}{2(\alpha''\xi+\beta''\eta+\gamma'')K}\bigl[\{\cos(\vartheta-\vartheta_{0})-i\sin(\vartheta-\vartheta_{0})\}\{(\xi-\lambda)+i(\eta-\mu)\}+\{\cos(\vartheta-\vartheta_{0})+i\sin(\vartheta-\vartheta_{0})\}\{(\xi-\lambda)-i(\eta-\mu)\}\bigr]
$} \]
\[ =\dfrac{k}{(\alpha''\xi+\beta''\eta+\gamma'')K}\bigl[(\xi-\lambda)\cos(\vartheta-\vartheta_{0})+(\eta-\mu)\sin(\vartheta-\vartheta_{0})\bigr], \]
where $J,J'$ have been replaced by
\[ J=\dfrac{k}{K}\cos\vartheta_{0},\quad J'=-\dfrac{k}{K}\sin\vartheta_{0}. \]

And putting besides
\[ \vartheta-\vartheta_{0}=\theta, \]
we have more simply
\[ (x-l)\cos\vartheta+(y-m)\sin\vartheta=\dfrac{k}{(\alpha''\xi+\beta''\eta+\gamma'')K}\bigl[(\xi-\lambda)\cos\theta+(\eta-\mu)\sin\theta\bigr], \]
whence
\[ (x-l)\cos\vartheta+(y-m)\sin\vartheta-\varpi \]
\[ =\dfrac{k}{(\alpha''\xi+\beta''\eta+\gamma'')K}\left\{(\xi-\lambda)\cos\theta+(\eta-\mu)\sin\theta-\dfrac{\varpi K(\alpha''\xi+\beta''\eta+\gamma'')}{k}\right\}. \]
Write now
\[ \alpha''=K\cos j,\quad \beta''=K\sin j;\quad a=\tfrac{k}{K},\quad b=\tfrac{k}{K}(\alpha''\lambda+\beta''\mu+\gamma''); \]
then
\[ (x-l)\cos\vartheta+(y-m)\sin\vartheta-\varpi \]
\[ =\sqrt{(\alpha''^{2}+\beta''^{2})}\,a\,\bigl[(\xi-\lambda)\cos\phi+(\eta-\mu)\sin(\phi-\varpi)\bigr], \]
where it will be noticed that
\[ (x-l)\cos\vartheta+(y-m)\sin\vartheta-\varpi=0, \]
and
\[ (\xi-\lambda)\cos\phi+(\eta-\mu)\sin\phi-\varpi=0, \]

% ============================================================
% Page 445 (PDF 467)
% ============================================================
\begin{center}
\textsc{283]} \hfill ON A THEOREM RELATING TO HOMOGRAPHIC FIGURES. \hfill 445
\end{center}

are lines in the first figure and in the second figure, corresponding to each other
and at the same distance $\varpi$ from the points $S,S'$ respectively. Call these
lines $T$ and $T'$.

But the relations between $\theta,\phi,\varpi$ are given by
\[ \varpi\cos j=a\cos\theta-b\cos\phi, \]
\[ \varpi\sin j=a\sin\theta-b\sin\phi, \]
or, by changing the fixed axes from which $\theta,\phi$ are respectively measured,
\[ \varpi=a\cos\theta-b\cos\phi, \]
\[ 0=a\sin\theta-b\sin\phi. \]

Now in these equations $\varpi$ is the perpendicular distance from $S$ upon the line
$(x-l)\cos\vartheta+(y-m)\sin\vartheta-\varpi=0$ and $\theta$ (which only differs from
$\vartheta$ by a constant angle) is the inclination of this perpendicular to a certain
fixed line; $\varpi$ is also the perpendicular distance of the line
$(\xi-\lambda)\cos\phi+(\eta-\mu)\sin\phi-\varpi=0$ from the point $S'$, and $\phi$
is the inclination of this perpendicular to a fixed line. Eliminating successively
$\phi$ and $\theta$, we have
\[ \varpi=a\cos\theta-\sqrt{(b^{2}-a^{2}\sin^{2}\theta)}, \]
\[ \varpi=-b\cos\phi+\sqrt{(a^{2}-b^{2}\sin^{2}\phi)}, \]
or, as these equations may also be written,
\[ \varpi^{2}-2\varpi a\cos\theta+a^{2}-b^{2}=0, \]
\[ \varpi^{2}+2\varpi b\cos\phi+b^{2}-a^{2}=0. \]

Suppose $a>b$, the former equation shows that the line $T$ is a tangent to a certain
hyperbola, and the latter equation shows that the line $T'$ is a tangent to a certain
ellipse, and it is easily seen that, taking for the transverse axes the lines from
which the angles $\theta$ and $\phi$ are respectively measured, the equation of the
hyperbola is
\[ \dfrac{x^{2}}{b^{2}}-\dfrac{y^{2}}{a^{2}-b^{2}}=1, \]
and that of the ellipse is
\[ \dfrac{\xi^{2}}{a^{2}}+\dfrac{\eta^{2}}{a^{2}-b^{2}}=1, \]
which are the conics $\Sigma,\Sigma'$ referred to in the enunciation. And if the
second conic is superimposed upon the first in such manner that the coordinates
$\xi,\eta$ may belong to the same axes with $x,y$, then the two conics will have the
assumed relation, viz.\ the foci of either conic will coincide with the vertices of
the other conic.

\bigskip
% ============================================================
% Page 446 (PDF 468) -- Paper 284 begins
% ============================================================
\begin{center}
446 \hfill \hfill \textsc{[284}
\end{center}

\begin{center}
\textbf{\large 284.}\\[1ex]
\textsc{ON A NEW ANALYTICAL REPRESENTATION OF CURVES IN SPACE.}
\end{center}

\begin{center}
\textit{[From the Quarterly Journal of Pure and Applied Mathematics, vol. \textsc{iii}.\ (1860),
pp. 225--236.]}
\end{center}

\medskip

The ordinary analytical representation of a curve in space by the equations of
two surfaces passing through the curve is, even in the case where the curve is the
complete intersection of the two surfaces, inappropriate as involving the consideration
of surfaces which are extraneous to the curve; and the objection becomes more serious
when the curve is not the complete intersection of any two surfaces; for in this case
the curve can only be represented in conjunction with a curve or curves extraneous
to itself. A curve in space can be in a mode represented by means of the equation
of the developable surface having the curve for its edge of regression; but this
corresponds to the representation of a plane curve by means of its equation in line
coordinates; a representation which is very useful in addition to, but which is not
to be substituted for, the equation in point coordinates. It occurred to me some
years ago that it might be advantageous to represent a curve in space by means of the
cone passing through the curve and having for its vertex an arbitrary
point\footnote{See my paper ``On the cones which pass through a curve of the third
order in space,'' \textit{Phil.\ Mag.}, t.\ \textsc{xii}.\ (1856), p.\ 20, [200], where
a curve of the third order in space is in effect so represented.}; and although I
have not advanced beyond the first steps of the theory, the results which I have
obtained may, I think, be interesting to geometers. The conclusion is that a curve
in space may be represented by a homogeneous equation $V=0$ between six coordinates
$(p,q,r,s,t,u)$ such that $ps+qt+ru=0$; the function $V$ is moreover such that in
virtue of this relation between the coordinates, and of the equation $V=0$ itself,
we have
\[ d_{p}V\cdot d_{s}V+d_{q}V\cdot d_{t}V+d_{r}V\cdot d_{u}V=0, \]

% ============================================================
% Page 447 (PDF 469)
% ============================================================
\begin{center}
\textsc{284]} \hfill ON A NEW ANALYTICAL REPRESENTATION OF CURVES IN SPACE. \hfill 447
\end{center}

or what is the same thing we have identically
\[ d_{p}V\cdot d_{s}V+d_{q}V\cdot d_{t}V+d_{r}V\cdot d_{u}V=LV+M(ps+qt+ru), \]
where $L$ and $M$ are functions of $p,q,r,s,t,u$. But the converse proposition, viz.\
any equation whatever $V=0$, where $V$ satisfies the condition just referred to
represents a curve in space, is not true; it would obviously be an important point
in the theory to ascertain what further conditions must be satisfied by the function
$V$.

The establishment of the foregoing results is very easy; in fact if $x,y,z,w$ are
current coordinates of the ordinary kind (point-coordinates) and $\alpha,\beta,\gamma,\delta$
the coordinates of an arbitrary point; the equation of any cone whatever having for
its vertex the point $(\alpha,\beta,\gamma,\delta)$ may be represented by a
homogeneous equation between the six determinants of the matrix
\[ \begin{pmatrix} x & y & z & w \\ \alpha & \beta & \gamma & \delta \end{pmatrix} \]
or if we write
\[ p=\gamma y-\beta z,\qquad s=\delta x-\alpha w, \]
\[ q=\alpha z-\gamma x,\qquad t=\delta y-\beta w, \]
\[ r=\beta x-\alpha y,\qquad u=\delta z-\gamma w, \]
values which it is well known give identically
\[ ps+qt+ru=0, \]
then the cone will be represented by a homogeneous equation
\[ V=0 \]
between the six coordinates $(p,q,r,s,t,u)$. It remains to find the conditions in
order that all the cones so represented, viz.\ the cones obtained by giving any
values whatever to the arbitrary quantities $\alpha,\beta,\gamma,\delta$ which enter
implicitly into the coordinates $p,q,r,s,t,u$, pass through one and the same curve;
for when this is the case, the equation $V=0$ may be properly considered as the
equation of the curve.

Assume then that all the cones pass through the same curve; if we give to one
of the arbitrary quantities $\alpha,\beta,\gamma,\delta$, say $\alpha$, the
infinitesimal variation $d\alpha$, then the function $V$ becomes $V+d_{\alpha}V\cdot d\alpha$,
and each of the equations $V=0$, $V+d_{\alpha}V\cdot d\alpha=0$ belongs to a cone
passing through the curve; the equation $d_{\alpha}V=0$ is therefore the equation of
a surface passing through the curve; and in like manner the four equations
\[ d_{\alpha}V=0,\quad d_{\beta}V=0,\quad d_{\gamma}V=0,\quad d_{\delta}V=0 \]
are each of them the equation of a surface passing through the curve, or these
equations must be simultaneously satisfied for all the points of the curve, they
must consequently reduce themselves to two independent relations. But $V$ is given as a

% ============================================================
% Page 448 (PDF 470)
% ============================================================
\begin{center}
448 \hfill ON A NEW ANALYTICAL REPRESENTATION OF CURVES IN SPACE. \hfill \textsc{[284}
\end{center}

function of $p,q,r,s,t,u$, through which quantities it is a function of $\alpha,\beta,\gamma,\delta$;
the last-mentioned four equations become therefore
\begin{align*}
0 &= \phantom{-}\, -d_{r}V\cdot y + d_{q}V\cdot z - d_{t}V\cdot w,\\
0 &= \phantom{-}\, \phantom{-}d_{r}V\cdot x - d_{p}V\cdot z - d_{u}V\cdot w,\\
0 &= -d_{q}V\cdot x + d_{p}V\cdot y - d_{u}V\cdot w,\\
0 &= \phantom{-}\, d_{s}V\cdot x + d_{t}V\cdot y + d_{u}V\cdot z,
\end{align*}
and from the first, second, and third equation, or any other combination of three
equations, we obtain at once the condition
\[ d_{p}V\cdot d_{s}V+d_{q}V\cdot d_{t}V+d_{r}V\cdot d_{u}V=0, \]
so that if this condition is satisfied, the four equations do in fact reduce
themselves to two independent equations; the condition in question is thus shown to
be necessary. But the condition only implies that all the cones having their vertices
in the neighbourhood of the point $(\alpha,\beta,\gamma,\delta)$ pass through one and
the same curve; this will be the case for instance for a series of cones all of them
circumscribed about one and the same surface, those having their vertices in the
neighbourhood of the point $(\alpha,\beta,\gamma,\delta)$ will all pass through the
curve of contact with the surface of the cone having for its vertex the point in
question. But the curve of contact is not a fixed curve for all positions of the
vertex, and the condition before referred to is consequently insufficient.

It may be noticed that the systems $p,q,r$ and $s,t,u$ are not similar to each
other and that the six coordinates cannot be in any way divided into two systems
which are similar to each other: the symmetry of the coordinates is in fact that of
the vertices (or sides) of a complete quadrilateral (or quadrangle); thus we may
divide the coordinates into two sets in a fourfold manner as follows:
\[ \begin{array}{l@{\quad}l}
u,t,p; & r,q,s,\\
s,q,u; & p,t,r,\\
r,t,s; & u,q,p,\\
p,q,r; & s,t,u,
\end{array} \]
where each left-hand set corresponds to three vertices forming a triangle and each
right-hand set to the remaining three vertices \textit{in lineo}. It may be noticed
also that if in the equation $V=0$ of any curve in space we substitute for
$p,q,r,s,t,u$, their values, and equate to zero the coefficients of the different
powers and products of $\alpha,\beta,\gamma,\delta$, each of the equations so
obtained will belong to a surface passing through the curve, and the entire system
of these equations will be equivalent to two relations only between the coordinates
$x,y,z,w$. But any two of these surfaces will not in general intersect only in the
curve, i.e.\ the curve will not be the complete intersection of any two of the
surfaces. It may be added that the equation of any other surface whatever through
the curve will be obtained by equating to zero a syzygetic function of the functions
which equated to zero give the surfaces first referred to.

% ============================================================
% Page 449 (PDF 471)
% ============================================================
\begin{center}
\textsc{284]} \hfill ON A NEW ANALYTICAL REPRESENTATION OF CURVES IN SPACE. \hfill 449
\end{center}

As an example of the theory, the equation, in the new coordinates, of a line in
space will be
\[ Ap+Bq+Cr+Fs+Gt+Hu=0, \]
where the constants are such that
\[ AF+BG+CH=0. \]

This may be verified as follows: suppose that the line in question is given as the
line of junction of the points $(\alpha',\beta',\gamma',\delta')$ and
$(\alpha'',\beta'',\gamma'',\delta'')$; then the cone through the line and the
arbitrary point $(\alpha,\beta,\gamma,\delta)$ is nothing else than the plane through
the three points; its equation is therefore
\[ \begin{vmatrix} x & y & z & w \\ \alpha & \beta & \gamma & \delta \\ \alpha' & \beta' & \gamma' & \delta' \\ \alpha'' & \beta'' & \gamma'' & \delta'' \end{vmatrix}=0, \]
which may be written
\[ Ap+Bq+Cr+Fs+Gt+Hu=0, \]
the values of $A$, \&c.\ being
\[ A=\alpha'\delta''-\alpha''\delta',\quad F=\beta'\gamma''-\beta''\gamma', \]
\[ B=\beta'\delta''-\beta''\delta',\quad G=\gamma'\alpha''-\gamma''\alpha', \]
\[ C=\gamma'\delta''-\gamma''\delta',\quad H=\alpha'\beta''-\alpha''\beta', \]
which in fact satisfy the relation
\[ AF+BG+CH=0. \]

So again if the line is given as the intersection of the planes
\[ ax+by+cz+dw=0, \]
\[ a'x+b'y+c'z+d'w=0, \]
then the equation of the cone (plane) through this line and the arbitrary point
$(\alpha,\beta,\gamma,\delta)$ is
\[ (ax+by+cz+dw)(a'\alpha+b'\beta+c'\gamma+d'\delta)-(a'x+b'y+c'z+d'w)(a\alpha+b\beta+c\gamma+d\delta)=0, \]
which is
\[ Ap+Bq+Cr+Fs+Gt+Hu=0, \]
the values of $A$, \&c.\ being
\[ A=bc'-b'c,\quad F=ad'-a'd, \]
\[ B=ca'-c'a,\quad G=bd'-b'd, \]
\[ C=ab'-a'b,\quad H=cd'-c'd, \]
which also satisfy the relation
\[ AF+BG+CH=0. \]

\medskip
\begin{center}C.\ IV.\hfill 57\end{center}

% ============================================================
% Page 450 (PDF 472)
% ============================================================
\begin{center}
450 \hfill ON A NEW ANALYTICAL REPRESENTATION OF CURVES IN SPACE. \hfill \textsc{[284}
\end{center}

I annex the following further investigation: let $p',q',r',s',t',u'$ and
$p'',q'',r'',s'',t'',u''$ be what $p,q,r,s,t,u$ become when $\alpha,\beta,\gamma,\delta$
are changed into $\alpha',\beta',\gamma',\delta'$ and $\alpha'',\beta'',\gamma'',\delta''$
respectively: the line represented by the equation
\[ Ap+Bq+Cr+Fs+Gt+Hu=0, \]
is also represented by the equations
\[ Ap'+Bq'+Cr'+Fs'+Gt'+Hu'=0, \]
\[ Ap''+Bq''+Cr''+Fs''+Gt''+Hu''=0, \]
which, reverting to the coordinates $x,y,z,w$, may be written
\[ \mathfrak{A}'x+\mathfrak{B}'y+\mathfrak{C}'z+\mathfrak{D}'w=0, \]
\[ \mathfrak{A}''x+\mathfrak{B}''y+\mathfrak{C}''z+\mathfrak{D}''w=0, \]
where
\[ \mathfrak{A}'=\phantom{-}C\beta'-B\gamma'+F\delta',\qquad \mathfrak{A}''=\phantom{-}C\beta''-B\gamma''+F\delta'', \]
\[ \mathfrak{B}'=-C\alpha'+A\gamma'+G\delta',\qquad \mathfrak{B}''=-C\alpha''+A\gamma''+G\delta'', \]
\[ \mathfrak{C}'=\phantom{-}B\alpha'-A\beta'+H\delta',\qquad \mathfrak{C}''=\phantom{-}B\alpha''-A\beta''+H\delta'', \]
\[ \mathfrak{D}'=-F\alpha'-G\beta'-H\gamma',\qquad \mathfrak{D}''=-F\alpha''-G\beta''-H\gamma'', \]
in which form the equations represent two planes each of them through the given
line: the equation of the cone (plane) through the given line and the arbitrary
point $(\alpha,\beta,\gamma,\delta)$ is
\[ (\mathfrak{A}'x+\mathfrak{B}'y+\mathfrak{C}'z+\mathfrak{D}'w)(\mathfrak{A}''\alpha+\mathfrak{B}''\beta+\mathfrak{C}''\gamma+\mathfrak{D}''\delta) \]
\[ -(\mathfrak{A}'\alpha+\mathfrak{B}'\beta+\mathfrak{C}'\gamma+\mathfrak{D}'\delta)(\mathfrak{A}''x+\mathfrak{B}''y+\mathfrak{C}''z+\mathfrak{D}''w)=0, \]
or, developing,
\[ \resizebox{\textwidth}{!}{$\displaystyle
(\mathfrak{B}'\mathfrak{C}''-\mathfrak{B}''\mathfrak{C}')p+(\mathfrak{C}'\mathfrak{A}''-\mathfrak{C}''\mathfrak{A}')q+(\mathfrak{A}'\mathfrak{B}''-\mathfrak{A}''\mathfrak{B}')r+(\mathfrak{A}'\mathfrak{D}''-\mathfrak{A}''\mathfrak{D}')s+(\mathfrak{B}'\mathfrak{D}''-\mathfrak{B}''\mathfrak{D}')t+(\mathfrak{C}'\mathfrak{D}''-\mathfrak{C}''\mathfrak{D}')u=0.
$} \]
Now
\[ \mathfrak{B}'\mathfrak{C}''-\mathfrak{B}''\mathfrak{C}'=(-C\alpha'+A\gamma'+G\delta')(B\alpha''-A\beta''+H\delta'') \]
\[ -(-C\alpha''+A\gamma''+G\delta'')(B\alpha'-A\beta'+H\delta'), \]
which putting
\[ \beta'\gamma''-\beta''\gamma'=p_{1},\quad \alpha'\delta''-\alpha''\delta'=s_{1}, \]
\[ \gamma'\alpha''-\gamma''\alpha'=q_{1},\quad \beta'\delta''-\beta''\delta'=t_{1}, \]
\[ \alpha'\beta''-\alpha''\beta'=r_{1},\quad \gamma'\delta''-\gamma''\delta'=u_{1}, \]
become
\[ A^{2}p_{1}+ABq_{1}+ACr_{1}-(BG+CH)s_{1}+AGt_{1}+AHu_{1}, \]
or as it may be written
\[ A(Ap_{1}+Bq_{1}+Cr_{1}+Fs_{1}+Gt_{1}+Hu_{1})-(AF+BG+CH)s_{1}. \]
Instead of writing at once $AF+BG+CH=0$, I denote it by $\nabla$ and I write also
\[ Ap_{1}+Bq_{1}+Cr_{1}+Fs_{1}+Gt_{1}+Hu_{1}=\Omega. \]

% ============================================================
% Page 451 (PDF 473)
% ============================================================
\begin{center}
\textsc{284]} \hfill ON A NEW ANALYTICAL REPRESENTATION OF CURVES IN SPACE. \hfill 451
\end{center}

The series of coefficients $\mathfrak{B}'\mathfrak{C}''-\mathfrak{B}''\mathfrak{C}'$, \&c.\ is
\[ A\Omega-\nabla s_{1},\quad B\Omega-\nabla t_{1},\quad C\Omega-\nabla u_{1}, \]
\[ F\Omega-\nabla p_{1},\quad G\Omega-\nabla q_{1},\quad H\Omega-\nabla r_{1}, \]
and the required equation, restoring for $\Omega,\nabla$ their values, is
\[ \resizebox{\textwidth}{!}{$\displaystyle
(Ap+Bq+Cr+Fs+Gt+Hu)(Ap_{1}+Bq_{1}+Cr_{1}+Fs_{1}+Gt_{1}+Hu_{1})
-(AF+BG+CH)(ps_{1}+p_{1}s+qt_{1}+q_{1}t+ru_{1}+r_{1}u)=0,
$} \]
which in virtue of $AF+BG+CH=0$ becomes
\[ (Ap+Bq+Cr+Fs+Gt+Hu)(Ap_{1}+Bq_{1}+Cr_{1}+Fs_{1}+Gt_{1}+Hu_{1})=0. \]

The equation
\[ Ap_{1}+Bq_{1}+Cr_{1}+Fs_{1}+Gt_{1}+Hu_{1}=0 \]
would imply that the two points $(\alpha',\beta',\gamma',\delta')$,
$(\alpha'',\beta'',\gamma'',\delta'')$ are in the same plane
\[ Ap+Bq+Cr+Fs+Gt+Hu=0, \]
(or, what is the same thing, that this line is intersected by the line through the
two points); in this exceptional case, the planes determined by the given line and
the two points respectively are one and the same plane, and they do not by their
intersection determine the given line. But in every other case the factor
$Ap_{1}+Bq_{1}+Cr_{1}+Fs_{1}+Gt_{1}+Hu_{1}$ is not equal to zero, and the foregoing
equation becomes
\[ Ap+Bq+Cr+Fs+Gt+Hu=0, \]
which is as it should be.

In what precedes it has been shown that the equation of the line through the
points $(\alpha',\beta',\gamma',\delta')$ and $(\alpha'',\beta'',\gamma'',\delta'')$ is
\[ Ap+Bq+Cr+Fs+Gt+Hu=0, \]
where
\[ A=\alpha'\delta''-\alpha''\delta',\quad F=\beta'\gamma''-\beta''\gamma', \]
\[ B=\beta'\delta''-\beta''\delta',\quad G=\gamma'\alpha''-\gamma''\alpha', \]
\[ C=\gamma'\delta''-\gamma''\delta',\quad H=\alpha'\beta''-\alpha''\beta', \]
and that the equation of the line of intersection of the planes $ax+by+cz+dw=0$,
$a'x+b'y+c'z+d'w=0$ is
\[ Ap+Bq+Cr+Fs+Gt+Hu=0, \]
where
\[ A=bc'-b'c,\quad F=ad'-a'd, \]
\[ B=ca'-c'a,\quad G=bd'-b'd, \]
\[ C=ab'-a'b,\quad H=cd'-c'd. \]

\medskip
\begin{center}57--2\end{center}

% ============================================================
% Page 452 (PDF 474)
% ============================================================
\begin{center}
452 \hfill ON A NEW ANALYTICAL REPRESENTATION OF CURVES IN SPACE. \hfill \textsc{[284}
\end{center}

It may be added that the condition of intersection of the two lines
\[ Ap+Bq+Cr+Fs+Gt+Hu=0, \]
\[ A'p+B'q+C'r+F's+G't+H'u=0, \]
is
\[ AF'+A'F+BG'+B'G+CH'+C'H=0, \]
and any other problems in relation to the line, for instance to find the condition
that a line may pass through a given point or lie in a given plane, \&c.\ may also be
solved by means of the new coordinates.

A curve of the second order in space is either a pair of lines or a plane conic,
each of which species may degenerate into a pair of intersecting lines. If we write
\[ W=Ap+Bq+Cr+Fs+Gt+Hu, \]
\[ W'=A'p+B'q+C'r+F's+G't+H'u, \]
then the equation of the pair of lines will be $V=WW'=0$, and it is worth while to
show that this value of $V$ satisfies the fundamental equation
$d_{p}V\cdot d_{s}V+d_{q}V\cdot d_{t}V+d_{r}V\cdot d_{u}V=0$; we have in fact
$d_{p}V=AW'+A'W$, \&c.\ and thence the left-hand side is
\[ \resizebox{\textwidth}{!}{$\displaystyle
W'^{2}(AF+BG+CH)+WW'(AF'+A'F+BG'+B'G+CH'+C'H)+W^{2}(A'F'+B'G'+C'H'),
$} \]
which in fact vanishes in virtue of the relations
\[ AF+BG+CH=0,\quad A'F'+B'G'+C'H'=0,\quad WW'=0. \]
The like proof applies to any curve made up of two or more curves.

Consider in the next place the plane conic given by the equations
\[ ax+by+cz+dw=0, \]
\[ x^{2}+y^{2}+z^{2}+w^{2}=0. \]
To find the equation of the cone having for its vertex the point $(\alpha,\beta,\gamma,\delta)$
and passing through the conic we must, according to Joachimsthal's method, substitute
in the two equations for $x,y,z,w$ the values $\lambda x+\mu\alpha$, $\lambda y+\mu\beta$,
$\lambda z+\mu\gamma$, $\lambda w+\mu\delta$, and from the resulting equations
eliminate $\lambda,\mu$. The two equations become
\[ \lambda^{2}(x^{2}+y^{2}+z^{2}+w^{2})+2\lambda\mu(\alpha x+\beta y+\gamma z+\delta w)+\mu^{2}(\alpha^{2}+\beta^{2}+\gamma^{2}+\delta^{2})=0, \]
\[ \lambda(ax+by+cz+dw)+\mu(a\alpha+b\beta+c\gamma+d\delta)=0, \]
and thence eliminating $\lambda,\mu$ we find
\[ (a\alpha+b\beta+c\gamma+d\delta)^{2}(x^{2}+y^{2}+z^{2}+w^{2}) \]
\[ -2(a\alpha+b\beta+c\gamma+d\delta)(ax+by+cz+dw)(\alpha x+\beta y+\gamma z+\delta w) \]
\[ +(ax+by+cz+dw)^{2}(\alpha^{2}+\beta^{2}+\gamma^{2}+\delta^{2})=0, \]

% ============================================================
% Page 453 (PDF 475)
% ============================================================
\begin{center}
\textsc{284]} \hfill ON A NEW ANALYTICAL REPRESENTATION OF CURVES IN SPACE. \hfill 453
\end{center}

where the left-hand side is as it should be a function of $p,q,r,s,t,u$; the equation
may in fact be written
\begin{align*}
&a^{2}(\phantom{p^{2}+}q^{2}+r^{2}+s^{2})
+b^{2}(p^{2}\phantom{+q^{2}}+r^{2}+t^{2})
+c^{2}(p^{2}+q^{2}\phantom{+r^{2}}+u^{2})
+d^{2}(p^{2}+q^{2}+r^{2}\phantom{+s^{2}})\\
&\quad +2bc(-qr+tu)+2ca(-rp+us)+2ab(-pq+st)\\
&\quad +2ad(qu-rt)+2bd(rs-pu)+2cd(pt-qs)=0
\end{align*}
which treated as a quadric in $(a,b,c,d)$ may be written
\[ \resizebox{\textwidth}{!}{$\displaystyle
\begin{pmatrix}
q^{2}+r^{2}+s^{2} & -pq+st & -rp+us & qu-rt \\
-pq+st & p^{2}+r^{2}+t^{2} & -qr+tu & rs-pu \\
-rp+us & -qr+tu & p^{2}+q^{2}+u^{2} & pt-qs \\
qu-rt & rs-pu & pt-qs & p^{2}+q^{2}+r^{2}
\end{pmatrix}(a,b,c,d)^{2}=0,
$} \]
or treated as a quadric in $(p,q,r,s,t,u)$, in the form
\[ \resizebox{\textwidth}{!}{$\displaystyle
\begin{pmatrix}
b^{2}+c^{2}+d^{2} & -ab & -ac & \cdot & -cd & bd \\
-ab & c^{2}+a^{2}+d^{2} & -bc & cd & \cdot & -ad \\
-ac & -bc & a^{2}+b^{2}+d^{2} & -bd & ad & \cdot \\
\cdot & cd & -bd & a^{2}+b^{2} & ab & ac \\
-cd & \cdot & ad & ab & a^{2}+c^{2} & bc \\
bd & -ad & \cdot & ac & bc & a^{2}+d^{2}
\end{pmatrix}(p,q,r,s,t,u)^{2}=0.
$} \]

Or again, in a form which is one of a system of four forms,
\begin{align*}
&(b^{2}+c^{2},\, c^{2}+a^{2},\, a^{2}+b^{2},\, -bc,\, -ca,\, -ab)(p,q,r)^{2}\\
&\quad +(b^{2}+d^{2},\, c^{2}+d^{2},\, c^{2}+d^{2},\, bc,\, ca,\, ab)(s,t,u)^{2}\\
&\quad +2\begin{pmatrix} \cdot & -cd & bd \\ cd & \cdot & -ad \\ -bd & ad & \cdot \end{pmatrix}\!(p,q,r)(s,t,u)=0.
\end{align*}

\end{document}
